\section{Introduction}

The innate ability of a relational program to run in multiple modes reduces the complicated task of \emph{finding} a solution to the much simpler task of \emph{checking} that a candidate satisfies the specification~\cite{lozov2019relational}. 
For example, a program which verifies whether a given sequence of vertices forms a path in a graph can enumerate paths. 
Likewise, a relation that checks whether some variable assignment satisfies a propositional formula can be run in reverse to search for satisfying assignments. 
A notable instance of this verifier-to-solver approach enables program synthesis: by running a relational interpreter backwards, we can generate programs whose evaluation yields a given value. 

One downside to this approach is its high overhead resulting in low performance. 
A relational interpreter should be carefully crafted to synthesize non-trivial programs within reasonable time constraints. 
Moreover, the implementation of \mk itself has to be heavily optimized which undermines its reputation for being easy to embed in any general-purpose language. 

The \mk community has devoted a lot of effort over the years to make program synthesis a reality.  
Among such efforts is the functional conversion~\cite{engel2023translation} whose purpose is to translate a relation with a given direction into a functional program. 
As a result, much of the overhead of relational programming is neutralized, especially if paired with advanced program transformation techniques such as specialization~\cite{verbitskaia2021deduction}. 

While the implementation of the functional conversion described in~\cite{engel2023translation} improves execution time, it exhibits several limitations. 
Firstly, because it relies on an untyped deep embedding of \mk in \haskell, it fails to take advantage of the \haskell type system, which leads to a loss of static guarantees as well as forces the user to painstakingly thread generator parameters. 
Secondly, the absence of automatic determinism analysis leaves some optimization opportunities untapped. 
Finally, the deep embedded representation hinders composition and extension---qualities that are among \mk's greater strengths. 

In this paper, we describe a typed, tagless-final~\cite{carette2009finally} embedding of \mk in \haskell. 
Along with automatic determinism analysis, it refines the earlier functional conversion, enhancing both performance and developer experience.
In addition to this, the embedding opens the door to composable and extensible \mk implementations in statically typed functional languages.
